% !TEX root = ../main.tex
\documentclass[../main.tex]{subfiles}

\begin{document}

%========================================================================================
% INTERMEZZO: UNIVERSITIES AS INTERSTITIAL ACTORS
%========================================================================================

\intermezzo{Universities as Interstitial Actors}

\begin{center}
\begin{tikzpicture}
    \fill[CoverGold] (-2,0) rectangle (-1.7,0.08);
    \fill[AccentPurple] (-1.5,0) rectangle (-1.2,0.08);
    \fill[AccentCyan] (-1,0) rectangle (-0.7,0.08);
    \fill[PandaBamboo] (-0.5,0) rectangle (-0.2,0.08);
    \fill[PandaRose] (0,0) rectangle (0.3,0.08);
\end{tikzpicture}
\end{center}

\vspace{1cm}

\begin{center}
\textit{\large Universities have always been interstitial.\\[0.3em]
They sit between knowledge production and knowledge application,\\[0.3em]
between state and market, between tradition and innovation.\\[0.3em]
What does the world's largest encyclopedia reveal about how\\[0.3em]
7,249 European universities narrate this interstitiality?}
\end{center}

\vspace{1.5cm}

\begin{chapteroverview}
This intermezzo applies text mining to a corpus of 7,249 English-language
Wikipedia articles on European universities across 38 countries (6.1 million words).
Using BERTopic for thematic discovery, named entity recognition for institutional
mapping, and event extraction for temporal analysis, we examine how Wikipedia's
editorial community describes the interstitial character of European higher
education. Three analyses structure the intermezzo: (1) a topic landscape
revealing the discursive terrain of European university descriptions;
(2) a named entity network showing which organizations, people, and places
co-occur across university articles; and (3) an event timeline extracting
founding dates, mergers, and reforms that mark institutional evolution. We
conclude by positioning the EIT partner universities within this broader
landscape, connecting back to the knowledge triangle analysis in
\cref{ch:interstitial-pluralism}.
\end{chapteroverview}


%===============================================================================
\section*{Why Wikipedia? Institutional Description at Scale}
\labsec{wiki-uni:rationale}
%===============================================================================

Universities are paradigmatic interstitial actors. They educate (producing human
capital), research (producing knowledge), and increasingly entrepreneurialize
(producing innovation and startups). This triple mandate---the ``knowledge
triangle'' that structures the EIT's design---predates the EIT by centuries.
Every university navigates it differently.

Wikipedia offers a unique vantage point for studying how this interstitiality
is narrated at scale.\sidenote{\textbf{Wikipedia as institutional
archive}---unlike websites (strategic self-presentation) or rankings (metric
reduction), Wikipedia articles are collaboratively authored third-party
descriptions that integrate historical, structural, and reputational
information.} Its articles are:

\begin{itemize}
    \item \textbf{Third-party authored}: not self-promotional, but written by
    external editors with encyclopedic intent
    \item \textbf{Structurally consistent}: following Wikipedia's manual of
    style for educational institutions, articles share comparable sections
    (history, academics, research, notable alumni)
    \item \textbf{Massively parallel}: 7,249 articles across 38 countries,
    enabling cross-national comparison at a scale impossible through manual
    case study
    \item \textbf{Temporally layered}: articles accumulate institutional
    history, from medieval founding to contemporary innovation programs
\end{itemize}

The corpus comprises 3,377 articles with 500+ words
(the analytically usable subset), totaling 6.1 million words,
with a median article length of 452 words.\sidenote{Articles were retrieved
programmatically via the Wikipedia API (\texttt{wikipedia-api} Python package),
traversing ``Category:Universities in [Country]'' trees for 38 European
countries. Filtering removed non-university articles (student unions,
libraries, sports teams) using title and summary heuristics.}


%===============================================================================
\section*{Topic Landscape: What Wikipedia Says About Universities}
\labsec{wiki-uni:topics}
%===============================================================================

We apply BERTopic \parencite{grootendorst2022bertopic} to discover thematic
structure in the corpus.\sidenote{\textbf{BERTopic pipeline}---documents
embedded with \texttt{all-MiniLM-L6-v2}, reduced via UMAP, clustered with
HDBSCAN (\texttt{min\_topic\_size=10}), and topic representations refined
using c-TF-IDF with English stop-word removal and bigram support.}
The model discovers 29 topics from the 3,377 substantial articles, assigning
73.6\% of documents to a topic (the remaining 26.4\% are outliers too
heterogeneous for any cluster).

The topic landscape reveals three dominant structuring principles:

\begin{center}
\includegraphics[width=0.95\textwidth]{figures/wiki-topic-landscape.pdf}
\captionof{figure}{BERTopic topic landscape of 3,377 European university Wikipedia
articles. Topics are labeled by their dominant terms and colored by category:
national-linguistic (gold), institutional type (purple), and British cultural
(cyan). The 29 topics assign 73.6\% of articles; the remaining 26.4\% are
outliers.}
\labfig{wiki-topic-landscape}
\end{center}

\paragraph{National-linguistic clustering.}
The largest topics partition by country or language group. Eastern European
and Ukrainian universities form the largest cluster (460 articles), followed by
British campus universities (392), Oxford--Cambridge colleges (153), German
applied sciences (151), Nordic universities (149), Italian universities (137),
Iberian universities (130), French grandes écoles (129), Irish universities (103),
and Benelux universities (75). Wikipedia's English-language descriptions of
European universities are, first and foremost, organized by national higher
education systems---each with its own institutional vocabulary, degree structures,
and historical narratives.

\paragraph{Institutional type specialization.}
Beyond national clustering, several topics capture institutional types that
cut across borders: arts and music academies (88 articles), theatre and design
schools (55), business schools (42), military academies (14), and
European politics and international relations programs (10). These
cross-national topics suggest that certain institutional identities---the
conservatoire, the business school, the military academy---carry sufficient
distinctiveness to override national framing.

\paragraph{British cultural specificity.}
A striking finding is the prevalence of distinctly British cultural
topics: rowing and boat clubs (44 articles), student newspapers (21),
student radio stations (19), university choirs (18), University Challenge
appearances on the BBC (18), academic dress (15), and students' unions (31).
These topics reflect not just the UK's numerical dominance in the corpus
(2,504 of 7,249 articles) but also Wikipedia's editorial community's
disproportionate attention to British university culture---a form of
encyclopedic Anglocentrism that shapes how the ``typical university'' is
described in the English-language Wikipedia.

\subsection*{EIT Partner Universities in the Topic Landscape}

Of the 27 EIT partner universities identified in the corpus, their topic
distribution is revealing: 6 fall in the Nordic cluster, 5 in the Iberian
cluster, 4 in the German applied sciences cluster, 3 each in the Eastern
European and British campus clusters, 2 each in the Oxford--Cambridge,
Italian, and Benelux clusters. No EIT partner universities appear in
the arts, business school, or cultural-specific topics. This dispersal
across primarily technical and research-oriented national clusters---rather
than concentration in any single topic---is consistent with the EIT's
selection criteria favoring technically oriented research universities
across diverse European countries.


%===============================================================================
\section*{Event Extraction: Institutional Histories}
\labsec{wiki-uni:events}
%===============================================================================

University Wikipedia articles are rich in temporal markers---founding dates,
mergers, name changes, reforms, and expansions. Using regex-based event
extraction on the first 5,000 characters of each article, we extract
4,189 events from the corpus: 2,601 founding dates, 1,353 institutional
renames, and 235 mergers.\sidenote{\textbf{Event extraction}---pattern matching
for ``founded/established/created in [YEAR]'', ``merged/amalgamated with'',
and ``renamed/became known as'' constructions. Founding years are validated
to the range 800--2025. Cross-validation against 24 well-known universities
yields 79\% accuracy (71\% exact match, 8\% within $\pm$5 years); errors
arise primarily from regex matching earlier dates mentioned in historical
context rather than the actual founding year.}

\paragraph{The long arc of European higher education.}
The extracted founding dates span nearly a millennium, from 1088 (the
University of Bologna, conventionally Europe's oldest) to 2025. The
distribution reveals clear historical waves:

\begin{itemize}
    \item \textbf{Medieval foundations} (1088--1499): 126 universities,
    concentrated in Italy, France, Spain, and the UK---the first European
    universities emerging from cathedral schools and monastic traditions
    \item \textbf{Early modern expansion} (1500--1799): 268 universities,
    driven by the Reformation (Protestant universities), Counter-Reformation
    (Jesuit colleges), and Enlightenment-era state building
    \item \textbf{Nineteenth-century massification} (1800--1899): 505
    universities, reflecting industrialization, the Humboldtian research
    university model, and the nation-state's need for trained professionals
    \item \textbf{Twentieth-century explosion} (1900--1999): 1,204
    universities---nearly half of all extracted founding dates---driven by
    post-war expansion, the welfare state, and the democratization of
    higher education
    \item \textbf{Twenty-first-century continuation} (2000--2025): 498
    universities, many in post-socialist countries, Turkey, and as mergers
    of existing institutions
\end{itemize}

The median founding year is 1953, placing the ``typical'' European university
squarely in the post-war expansion era.

\begin{center}
\includegraphics[width=0.95\textwidth]{figures/wiki-founding-timeline.pdf}
\captionof{figure}{Distribution of university founding dates extracted from
Wikipedia articles (25-year bins). Colors indicate historical eras: medieval
and early modern (gold), nineteenth century (gold, lighter), twentieth century
(purple), and twenty-first century (cyan). The dashed line marks the median
founding year (1953).}
\labfig{wiki-founding-timeline}
\end{center}

\paragraph{National temporal signatures.}
Countries exhibit distinctive temporal signatures that reflect their
political histories.\sidenote{Median founding year by country (selected):
Albania 2004, Italy 1935, UK 1928, France 1941, Germany 1968,
Ukraine 1930, Turkey 1987, Luxembourg 2003.} The UK (median 1928) and
Italy (1935) have older median founding dates, reflecting deep
institutional histories. Germany (1968) and Austria (1967) show
the post-war Fachhochschule expansion. Post-socialist countries
cluster in the 1950s--1990s (Romania 1948, Czech Republic 1968,
Latvia 1993, Slovenia 1995). The youngest systems---Albania (2004),
Montenegro (2004), North Macedonia (2007)---reflect post-independence
institution building.

\begin{center}
\includegraphics[width=0.95\textwidth]{figures/wiki-country-founding.pdf}
\captionof{figure}{Median founding year by country (countries with $\geq 5$
extracted events). Vertical reference lines mark 1945 (post-war) and 1992
(polytechnic conversion / post-socialist transition). Countries range from
Italy and the UK (oldest median) to Albania and North Macedonia (youngest).}
\labfig{wiki-country-founding}
\end{center}

\paragraph{Mergers and renames as interstitial markers.}
The 235 mergers and 1,353 renames reveal institutions in transition.
Renaming is particularly common in post-socialist countries (where
universities shed Soviet-era names) and in the UK (where polytechnics
became universities after the 1992 Further and Higher Education Act).
Mergers accelerate in the 2000s and 2010s, driven by efficiency pressures,
the Bologna Process, and excellence initiatives---institutional
restructuring that creates the interstitial conditions the EIT later
exploits.

\begin{center}
\includegraphics[width=0.95\textwidth]{figures/wiki-event-decades.pdf}
\captionof{figure}{Institutional events by decade: founding dates, renames,
and mergers extracted from Wikipedia articles. Founding dates peak in the
1960s--1990s (post-war expansion); renames cluster in the 1990s--2000s
(post-socialist and polytechnic transitions); mergers accelerate after 2000.}
\labfig{wiki-event-decades}
\end{center}


%===============================================================================
\section*{Named Entity Network: Institutional Connections}
\labsec{wiki-uni:ner}
%===============================================================================

Named entity recognition (spaCy, \texttt{en\_core\_web\_sm}) extracts
397,252 entity mentions across the 3,377 articles, comprising 163,426
unique entities.\sidenote{\textbf{NER pipeline}---spaCy's small English model
applied to the first 10,000 characters of each article. Entity types retained:
ORG, GPE (geopolitical entities), PERSON, DATE.} Co-occurrence of
organizations within the same article creates an implicit institutional
network: 41,138 edges with weight $\geq 3$ (i.e., organizations mentioned
together in at least three different university articles).

\paragraph{Most-mentioned organizations.}
The top organizations reveal the institutional landmarks of European
higher education: Oxford (1,218 mentions), the University of Cambridge
(237), the University of London (235), BBC (229), UCL (228), and the
EU (213). The presence of the BBC reflects the UK corpus dominance and
British university culture (University Challenge, student media). The
EU's prominence signals the growing Europeanization of higher
education discourse.

\paragraph{Geographic anchors.}
The most-mentioned locations trace the urban geography of European
higher education: UK (1,568), Germany (1,277), London (1,233),
France (1,119), Cambridge (1,109). These are not merely country names
but reflect the gravitational centers around which university
descriptions orbit---London as the hub of the British system, Paris
for France, Berlin and Munich for Germany.

\paragraph{Co-occurrence network structure.}
The 41,138 co-occurrence edges create a dense network of implicit
institutional relationships. Universities that appear in the same
Wikipedia articles share contextual space---they are compared, grouped,
or historically connected by Wikipedia's editorial community. This
network, while not analyzed with full community detection here, provides
a complementary view to the website-based network in
\cref{ch:interstitial-pluralism}: where website text reveals how
organizations describe \emph{themselves}, the co-occurrence network
reveals how they are described \emph{in relation to each other} by an
external editorial community.


%===============================================================================
\section*{SeeAlsology: The Wikipedia Link Network}
\labsec{wiki-uni:seealsology}
%===============================================================================

Wikipedia's ``See also'' sections create an editorial link network---how
Wikipedia's community of editors conceptually connects universities to each
other and to related concepts. Inspired by the SeeAlsology
tool,\sidenote{\textbf{SeeAlsology}---a visual network exploration tool by
DensityDesign Lab that traces ``See also'' links on Wikipedia, revealing
conceptual proximity as perceived by Wikipedia editors
\parencite{ricci2014seealsology}. We replicate its logic programmatically via
the MediaWiki parse API.} we extract ``See also'' links from each of the 34
EIT partner university articles and follow them one hop out, building a
two-depth neighborhood network.

\begin{center}
\includegraphics[width=0.95\textwidth]{figures/wiki-seealsology-network.pdf}
\captionof{figure}{SeeAlsology network of EIT partner universities' ``See also''
links (two-hop neighborhood). Gold nodes are EIT seed universities; purple nodes
are bridge concepts. Node size reflects degree. Only one direct EIT-to-EIT link
exists (Chalmers $\rightarrow$ KTH). The Bologna Process (degree 31) and
national university lists serve as primary bridges.}
\labfig{wiki-seealsology}
\end{center}

The resulting network comprises 452 nodes and 578 edges. Of the 34 seed
universities, 27 have ``See also'' sections (the remaining 7 lack this
section entirely---a signal of thinner editorial attention). The network
reveals three patterns:

\paragraph{Sparse EIT interconnection.}
Only one direct ``See also'' link connects two EIT partner universities:
Chalmers University of Technology $\rightarrow$ KTH Royal Institute of
Technology (both Swedish). This near-absence of editorial cross-linking
confirms that Wikipedia's editors do not perceive EIT partners as a
coherent group---they are connected to their national higher education
contexts, not to each other. The EIT as institutional network is
\emph{invisible} to Wikipedia's editorial community.

\paragraph{The EIT as bridge concept.}
The ``European Institute of Innovation and Technology'' article itself
appears as a bridge node with degree 13 in the network---linked from
or to 13 other articles in the two-hop neighborhood. This is notable:
while the EIT partner universities are not linked to each other, they
are indirectly connected \emph{through} the EIT's own Wikipedia article
and through shared concepts like the Bologna Process (degree 31), lists
of national universities, and the concept of ``institutes of technology''
(degree 10).

\paragraph{National and thematic bridges.}
The highest-degree bridge concepts are national university lists
(``List of universities in the United Kingdom'' at degree 24,
``List of universities in Germany'' at 15), educational policy
(``Bologna Process'' at 31, ``Education in Germany'' at 13),
and cultural concepts (``Town and gown'' at 20). The concept
``Science and technology in Flanders'' (degree 13) and
``Science and technology in Switzerland'' (degree 10) reveal
regional innovation ecosystem framing. The presence of the
``Vlaams Instituut voor Biotechnologie'' (VIB, degree 18) as
a bridge node illustrates how research institutes mediate
between university articles---a structural analogue to the
knowledge triangle's research-education interface.


%===============================================================================
\section*{Connecting Back: Universities in the Knowledge Triangle}
\labsec{wiki-uni:connection}
%===============================================================================

This intermezzo provides the broader landscape for
\cref{ch:interstitial-pluralism}'s analysis of EIT Community positioning.
Where that chapter examines how 1,298 entities position themselves through
website text, this intermezzo shows how 7,249 universities are described
through the lens of Wikipedia's editorial community. Three bridging
observations emerge:

\begin{enumerate}
    \item \textbf{National systems dominate the topic landscape.}
    The 29 BERTopic topics organize primarily by country and language group,
    not by knowledge-triangle pole. Wikipedia does not describe universities
    primarily as ``research universities'' or ``teaching universities'' but as
    \emph{national} institutions embedded in specific higher education
    systems. This suggests that the knowledge triangle framing---central to
    the EIT's design---is an \emph{overlay} on more fundamental national
    organizing logics, not the natural vocabulary through which universities
    are described.

    \item \textbf{The event timeline reveals waves of interstitialization.}
    The founding date distribution shows that most European universities
    (median 1953) are relatively young institutions. The post-war expansion,
    the polytechnic-to-university conversions of the 1990s, and the merger
    waves of the 2000s--2010s represent successive rounds of institutional
    reconfiguration---each creating new interstitial positions between
    education, research, and business. The EIT arrives not as a creator of
    interstitiality but as a \emph{formalizer} of conditions already
    produced by decades of higher education reform.

    \item \textbf{EIT partners are dispersed, not clustered.}
    The 27 EIT partner universities in the corpus spread across 8 of the 29
    topics, concentrated in national-technical clusters rather than forming
    a distinctive topic of their own. This dispersal is analytically
    significant: it confirms that EIT membership does not correspond to a
    distinctive Wikipedia discursive profile but rather selects from the
    technically oriented segments of diverse national systems.
\end{enumerate}

\closeintermezzo

\end{document}
